\chapter{Autenticazione, SSL/TLS e sicurezza operativa}

Ricordiamo le quattro proprietà di una comunicazione sicura: \textbf{confidenzialità} (contro lo sniffing), \textbf{integrità del
messaggio}, \textbf{autenticazione degli estremi} (contro lo spoofing) e \textbf{sicurezza operativa}. Le prime tre sono realizzate via
software; l'ultima si basa tipicamente su hardware dedicato. In questo capitolo completiamo l'autenticazione, vediamo come tutto si
combina in SSL/TLS e chiudiamo con firewall e IDS.

% ══════════════════════════════════════════════════════════════
\section{L'autenticazione degli estremi}
% ══════════════════════════════════════════════════════════════

\dfn{Autenticazione degli estremi}{
    L'\textbf{autenticazione degli estremi} (\emph{end-point authentication}) è il processo con cui un'entità \textbf{dimostra la propria
    identità} a un'altra entità attraverso una rete.
}

\begin{itemize}
    \item Tipicamente un \textbf{protocollo di autenticazione} (AP, \emph{Authentication Protocol}) viene eseguito \textbf{prima} che le due
          parti eseguano qualche altro protocollo: serve a stabilire le identità di entrambe prima di cominciare a lavorare.
    \item Esempi di protocolli che ne hanno bisogno: un trasferimento affidabile (basato su TCP), uno scambio di informazioni di routing (DV o
          LS), un protocollo di posta (SMTP), \dots
    \item Per l'autenticazione ci si può appoggiare alle \textbf{CA} viste nel capitolo precedente, ma non basta.
\end{itemize}

% ══════════════════════════════════════════════════════════════
\section{Tentativi progressivi di autenticazione}
% ══════════════════════════════════════════════════════════════

Supponiamo che l'host A debba autenticarsi presso B prima di lavorare insieme. Costruiamo un protocollo per tentativi, correggendo a ogni
passo il difetto del precedente (l'approccio del Kurose, con i protocolli ap1.0, ap2.0, \dots).

\subsection{Tentativo 1: controllare l'indirizzo IP}

Se A e B hanno già comunicato e l'indirizzo IP di A è sempre lo stesso, B potrebbe autenticare A \textbf{controllando l'indirizzo IP} nel
datagramma: se è quello noto, si assume che a inviare sia A.

\warning{Basta falsificare l'indirizzo}{
    Questa tecnica funziona contro intrusi ``ingenui'' che cercano di spacciarsi per A, ma non basta: un intruso ``meno ingenuo'' può
    \textbf{falsificare l'indirizzo} nel pacchetto (spoofing).
}

Un modo di evitare lo spoofing è sfruttare i \textbf{router}: un router può essere configurato per inoltrare solo datagrammi legittimi, cioè con
indirizzo sorgente che appartiene davvero agli host dietro di lui.

\begin{itemize}
    \item Il router deve conoscere gli IP dei propri host (altrimenti gli host sarebbero irraggiungibili).
    \item Può quindi rilevare i pacchetti il cui indirizzo non è coerente con la sorgente, e scartarli prima che facciano danni.
\end{itemize}

Purtroppo questa capacità \textbf{non è universalmente adottata} né imposta: bisogna assumere che lo spoofing sia possibile.

\subsection{Tentativo 2: una password segreta}

Un altro approccio, il più comune, è usare una \textbf{password segreta}: la password funziona come un segreto condiviso tra chi autentica e chi
viene autenticato. Gmail, Facebook, Telnet, FTP e moltissimi altri servizi usano l'autenticazione con password. Lo stesso segreto può servire
anche al controllo di integrità.

\warning{La password si può intercettare}{
    Il primo problema è che l'intruso può \textbf{intercettare la password} (eavesdropping) dai messaggi, e usarla insieme a un IP falsificato. Con
    Telnet, che trasmette tutto in chiaro, basta uno sniffer.
}

\subsection{Tentativo 3: una password cifrata}

L'intercettazione si evita \textbf{cifrando la password}, così che diventi illeggibile per l'intruso. Se si usa una cifratura simmetrica non serve
nemmeno una password aggiuntiva: la chiave condivisa stessa, che deve essere un segreto tra i due host, può fare da password.

\warning{L'attacco a ripetizione (playback)}{
    La cifratura è una buona cosa, ma non basta: un intruso ``molto furbo'' può infilarsi nella comunicazione fingendosi A con un
    \textbf{attacco a ripetizione} (\emph{playback} o \emph{replay attack}).
    \begin{itemize}
      \item L'intruso sniffa i pacchetti da A a B cercando di ottenere la versione \textbf{cifrata} della password di A.
      \item Se ci riesce, può \textbf{ripetere} la password cifrata a B in una nuova sessione, usandola \textbf{senza nemmeno capirla}.
      \item A differenza dei MitM classici (in tempo reale), sniffing e ripetizione avvengono in momenti separati.
    \end{itemize}
    Il problema è che B non riesce a capire se A è ``\textbf{vivo}'', cioè se c'è una sessione attiva in corso tra il vero A e B.
}

\begin{figure}[H]
\centering
\begin{tikzpicture}[yscale=0.8, every node/.style={font=\scriptsize\sffamily}]
  \timeline{a}{0}{A}{4.5}
  \timeline{t}{4}{Trudy}{4.5}
  \timeline{b}{8}{B}{4.5}
  \draw[msg, netblue] (0,-0.4) -- node[above, pos=0.3] {``sono A'', $K_{AB}$(password)} (8,-0.4);
  \fill[netred] (4,-0.4) circle (2.5pt); \node[text=netred, anchor=west] at (4.1,-0.8) {registra};
  \node[text=black!50] at (4,-1.6) {\dots giorni dopo \dots};
  \draw[msg, netred] (4,-2.6) -- node[above] {``sono A'', $K_{AB}$(password)} (8,-2.6);
  \node[anchor=west, text width=2.8cm, align=left] at (8.2,-2.6) {B decifra: la password è corretta!};
\end{tikzpicture}
\caption{L'attacco playback: Trudy non conosce la password, ma riusa il messaggio cifrato.}
\end{figure}

Il problema somiglia all'apertura di una connessione TCP, in cui con il three-way handshake i due host si assicurano di essere entrambi vivi. Nel
three-way handshake gli host usano \textbf{numeri di sequenza iniziali casuali}, perché ritrasmissioni di vecchie connessioni non vengano
scambiate per SYN o SYNACK nuovi. Un'idea simile si può usare per l'autenticazione.

% ══════════════════════════════════════════════════════════════
\section{La soluzione: il nonce}
% ══════════════════════════════════════════════════════════════

\dfn{Nonce}{
    Un \textbf{nonce} è un numero casuale (o pseudocasuale) che un protocollo usa \textbf{una sola volta nella vita} (\emph{number used once}).
}

\begin{figure}[H]
\centering
\begin{tikzpicture}[yscale=0.85, every node/.style={font=\scriptsize\sffamily}]
  \timeline{a}{0}{Host A}{4.6}
  \timeline{b}{7}{Host B}{4.6}
  \node[anchor=east] at (-0.2,-0.2) {chiave $K_{AB}$}; \node[anchor=west] at (7.2,-0.2) {chiave $K_{AB}$};
  \draw[msg] (0,-0.5) -- node[above] {1. ``sono A''} (7,-1.0);
  \node[draw, fill=llink!60, anchor=west] at (7.2,-1.3) {2. crea un nonce $R$ usato una sola volta};
  \draw[msg, netorange] (7,-1.6) -- node[above] {$R$} (0,-2.1);
  \draw[msg, netblue] (0,-2.6) -- node[above] {3. $K_{AB}(R + \text{password})$} (7,-3.1);
  \node[draw, fill=netgreen!20, anchor=west, text width=3.6cm, align=left] at (7.2,-3.6) {4. decifra: se il nonce è quello inviato, dall'altra parte c'è davvero A};
\end{tikzpicture}
\caption{Autenticazione con nonce e chiave simmetrica condivisa.}
\end{figure}

Supponendo che A e B condividano una chiave simmetrica, il nonce si usa così:

\begin{enumerate}
    \item A invia a B il messaggio ``sono A''.
    \item B sceglie un nonce $R$ e lo invia ad A.
    \item A cifra il nonce insieme alla password con la chiave simmetrica condivisa e rimanda il risultato a B.
    \item B decifra il messaggio: se il nonce decifrato coincide con quello che aveva inviato, A è autenticato.
\end{enumerate}

\begin{itemize}
    \item La password diventa una sorta di \textbf{password monouso}: ogni password cifrata è legata al proprio nonce, e non può essere riusata in altre
          sessioni (da qualcun altro).
    \item Se password e nonce arrivano correttamente a B, si può ragionevolmente assumere che dietro ci sia A, \textbf{vivo} in quel momento.
\end{itemize}

\info{Con la crittografia asimmetrica}{
    Lo stesso schema funziona con le chiavi pubbliche: B invia il nonce $R$, A lo cifra con la propria \textbf{chiave privata}, e B lo decifra con la
    chiave \textbf{pubblica} di A. Solo A può aver prodotto quel messaggio. Resta però il problema di essere sicuri che la chiave pubblica sia davvero
    di A: di nuovo servono i \textbf{certificati}.
}

% ══════════════════════════════════════════════════════════════
\section{SSL e TLS}
% ══════════════════════════════════════════════════════════════

Abbiamo visto che diverse tecniche crittografiche si possono usare e integrare per offrire confidenzialità, integrità e autenticazione alle
connessioni TCP.

\dfn{SSL/TLS}{
    Queste tecniche sono realizzate dal \textbf{Secure Socket Layer} (SSL), una versione ``potenziata'' di TCP che offre servizi di sicurezza.
    \begin{itemize}
      \item Una versione leggermente modificata di SSL 3 si chiama \textbf{Transport Layer Security} (TLS), ed è quella usata oggi.
      \item Esiste un equivalente per le comunicazioni basate su UDP, il \textbf{Datagram TLS} (DTLS).
    \end{itemize}
}

\begin{itemize}
    \item SSL è usato da molti protocolli applicativi: Web, posta, messaggistica istantanea, voice over IP, \dots
    \item Spesso la connessione tra il browser e un sito usa \textbf{HTTPS} (HTTP + SSL/TLS) invece di HTTP: il lucchetto nella barra degli indirizzi.
\end{itemize}

\begin{figure}[H]
\centering
\begin{tikzpicture}[every node/.style={font=\footnotesize\sffamily}]
  \begin{scope}
    \node[font=\bfseries\sffamily] at (0,2.3) {TCP normale};
    \node[lay app, minimum width=3cm] (a1) at (0,1.5) {Applicazione};
    \node[lay tra, minimum width=3cm] (t1) at (0,0.6) {TCP};
    \node[lay net, minimum width=3cm] (n1) at (0,-0.3) {IP};
    \node[lay link, minimum width=3cm] (l1) at (0,-1.2) {Accesso (link + fisico)};
    \draw[very thick, netblue] (-1.6,1.05) -- (1.6,1.05); \node[right, text=netblue] at (1.6,1.05) {socket TCP};
  \end{scope}
  \begin{scope}[xshift=7.5cm]
    \node[font=\bfseries\sffamily] at (0,2.3) {TCP con SSL};
    \node[lay app, minimum width=3cm] (a2) at (0,1.8) {Applicazione};
    \node[layer, fill=netred!30, minimum width=3cm] (s2) at (0,1.0) {sottolivello SSL};
    \node[lay tra, minimum width=3cm] (t2) at (0,0.2) {TCP};
    \node[lay net, minimum width=3cm] (n2) at (0,-0.6) {IP};
    \node[lay link, minimum width=3cm] (l2) at (0,-1.4) {Accesso (link + fisico)};
    \draw[very thick, netred] (-1.6,1.4) -- (1.6,1.4); \node[right, text=netred] at (1.6,1.4) {socket SSL};
    \draw[very thick, netblue] (-1.6,0.6) -- (1.6,0.6); \node[right, text=netblue] at (1.6,0.6) {socket TCP};
    \draw[decorate, decoration={brace, amplitude=5pt, mirror}] (-1.7,2.1) -- (-1.7,0.7) node[midway, left=6pt, align=right] {livello\\applicativo};
  \end{scope}
\end{tikzpicture}
\caption{SSL si inserisce tra l'applicazione e TCP: tecnicamente sta nel livello applicativo, ma allo sviluppatore appare come un trasporto sicuro.}
\end{figure}

\begin{itemize}
    \item Poiché SSL protegge TCP, può essere usato da \textbf{qualsiasi applicazione} che gira su TCP.
    \item SSL offre un'\textbf{API} semplice, con socket simili a quelli di TCP: un'applicazione che vuole usare SSL include le librerie SSL.
    \item Anche se SSL risiede tecnicamente nel livello applicativo, dal punto di vista dello sviluppatore sembra un protocollo di trasporto che offre i
          servizi di TCP arricchiti di sicurezza.
\end{itemize}

SSL ha tre fasi principali: \textbf{handshake}, \textbf{derivazione delle chiavi} e \textbf{trasferimento dati}.

% ══════════════════════════════════════════════════════════════
\section{SSL: l'handshake}
% ══════════════════════════════════════════════════════════════

Un client B vuole usare SSL per comunicare con un server A. Nella fase di handshake B deve:

\begin{itemize}
    \item \textbf{stabilire una connessione TCP} con A;
    \item \textbf{verificare che A sia davvero A} (e non un server fasullo);
    \item \textbf{inviare ad A una chiave segreta principale} (\emph{master secret}), il segreto condiviso.
\end{itemize}

\begin{figure}[H]
\centering
\begin{tikzpicture}[yscale=0.72, every node/.style={font=\scriptsize\sffamily}]
  \timeline{b}{0}{Client B}{8.4}
  \timeline{a}{7}{Server A}{8.4}
  \draw[msg, black!60] (0,-0.3) -- node[above] {TCP SYN} (7,-0.8);
  \draw[msg, black!60] (7,-1.0) -- node[above] {TCP SYNACK} (0,-1.5);
  \draw[msg, black!60] (0,-1.7) -- node[above] {TCP ACK} (7,-2.2);
  \draw[msg, netblue, thick] (0,-2.8) -- node[above] {SSL hello} (7,-3.3);
  \draw[msg, netgreen!60!black, thick] (7,-3.6) -- node[above] {certificato (con la chiave pubblica di A)} (0,-4.1);
  \node[draw, fill=llink!60, anchor=east, text width=2.8cm, align=right] at (-0.2,-4.8) {verifica il certificato con la CA; crea il Master Secret MS};
  \draw[msg, netred, thick] (0,-5.6) -- node[above] {$K_A^+$(MS): MS cifrato con la chiave pubblica di A} (7,-6.1);
  \node[draw, fill=llink!60, anchor=west, text width=2.8cm, align=left] at (7.2,-6.5) {decifra con la chiave privata e ottiene MS};
  \node[draw, fill=netgreen!20, text width=5cm, align=center] at (3.5,-7.6) {ora solo A e B conoscono MS};
\end{tikzpicture}
\caption{L'handshake SSL (semplificato).}
\end{figure}

\begin{enumerate}
    \item Stabilita la connessione TCP, B invia ad A un messaggio \textbf{hello}.
    \item Il server A risponde con il proprio \textbf{certificato}, che contiene la sua \textbf{chiave pubblica}.
    \item Poiché il certificato è stato certificato da una CA, il client B si fida che la chiave pubblica nel certificato appartenga ad A.
    \item B genera un \textbf{Master Secret} (MS), che userà solo per questa sessione SSL (è un nonce), lo cifra con la chiave pubblica di A e lo invia ad A.
    \item Il server A decifra il messaggio con la propria chiave privata e ottiene MS. Dopo questa fase, \textbf{solo B e A} conoscono il master secret della sessione.
\end{enumerate}

\info{Asimmetrica per iniziare, simmetrica per lavorare}{
    SSL usa la crittografia \textbf{asimmetrica} (lenta) solo per scambiare in sicurezza il segreto, e poi la crittografia \textbf{simmetrica} (veloce) per
    tutti i dati. È lo schema ibrido tipico di quasi tutti i protocolli sicuri.
}

% ══════════════════════════════════════════════════════════════
\section{SSL: la derivazione delle chiavi}
% ══════════════════════════════════════════════════════════════

Poiché MS è un segreto condiviso tra B e A, potrebbe essere usato come chiave simmetrica per tutta la cifratura e il controllo di integrità della
sessione. È però più sicuro che A e B usino \textbf{chiavi diverse} per la cifratura e per l'integrità, e per i due flussi di dati (da A a B e da B ad A).

Da MS si ricavano \textbf{quattro chiavi}:

\begin{table}[H]
\centering
\begin{tabular}{ll}
\toprule
\textbf{Chiave} & \textbf{Uso} \\
\midrule
$E_{B\rightarrow A}$ & chiave di \textbf{cifratura} di sessione per i dati da B ad A \\
$M_{B\rightarrow A}$ & chiave \textbf{MAC} di sessione per i dati da B ad A \\
$E_{A\rightarrow B}$ & chiave di \textbf{cifratura} di sessione per i dati da A a B \\
$M_{A\rightarrow B}$ & chiave \textbf{MAC} di sessione per i dati da A a B \\
\bottomrule
\end{tabular}
\caption{Le quattro chiavi di sessione derivate dal master secret (MAC = Message Authentication Code).}
\end{table}

% ══════════════════════════════════════════════════════════════
\section{SSL: il trasferimento dati}
% ══════════════════════════════════════════════════════════════

SSL protegge TCP spezzando il flusso di dati dell'applicazione in \textbf{record}. Per ciascun record:

\begin{enumerate}
    \item SSL accoda al record un \textbf{MAC}, per il controllo di integrità ($\text{record}_i + \text{MAC}$);
    \item il record così modificato viene \textbf{cifrato};
    \item il record cifrato viene passato a TCP per la trasmissione.
\end{enumerate}

\warning{E l'integrità dell'intero flusso?}{
    Così si garantisce l'integrità dei singoli record, ma non dell'intero flusso. Poiché è cifrato solo il payload, un MitM tra A e B può ancora
    danneggiare il flusso \textbf{invertendo o eliminando segmenti}. Per esempio può catturare due segmenti consecutivi, scambiarne l'ordine insieme ai
    numeri di sequenza TCP, e inviarli al ricevente, che non si accorge dell'inversione finché i dati non arrivano all'applicazione.
}

Per evitarlo SSL integra anche un \textbf{numero di sequenza} nel messaggio prima della cifratura: il vero messaggio cifrato è
$$\text{record}_i + \text{MAC} + \text{numeroDiSequenza}_i$$
dove il MAC è calcolato anche sul numero di sequenza. Conoscendo la procedura e la chiave MAC, il ricevente può certificare l'integrità e
l'\textbf{ordine} dei messaggi.

\begin{figure}[H]
\centering
\begin{tikzpicture}[every node/.style={font=\scriptsize\sffamily}]
  \node[field, fill=lapp!60, minimum width=6cm] (f) at (0,0) {flusso di dati dell'applicazione};
  \foreach \i/\x in {1/-2, 2/0, 3/2} \node[field, fill=lapp!60, minimum width=1.8cm] (r\i) at (\x,-1) {record \i};
  \draw[-{Stealth}] (0,-0.35) -- (0,-0.65);
  \node[field, fill=lapp!60, minimum width=1.8cm] (rr) at (-2,-2.2) {record 1};
  \node[field, fill=netred!25, minimum width=0.9cm, right=0pt of rr] (mac) {MAC};
  \node[field, fill=llink, minimum width=0.9cm, right=0pt of mac] (sq) {seq 1};
  \draw[decorate, decoration={brace, amplitude=5pt, mirror}] (rr.south west) -- (sq.south east) node[midway, below=6pt] {cifrato con $E_{B\rightarrow A}$ \quad (MAC calcolato con $M_{B\rightarrow A}$)};
  \node[field, fill=ltra, minimum width=1.2cm] (tcp) at (4.4,-2.2) {header TCP};
  \node[field, fill=black!15, minimum width=2.6cm, right=0pt of tcp] {payload cifrato};
  \draw[-{Stealth}, thick] (1.9,-2.2) -- (3.75,-2.2);
\end{tikzpicture}
\caption{Il trasferimento dati in SSL: record, MAC, numero di sequenza e cifratura, poi TCP.}
\end{figure}

% ══════════════════════════════════════════════════════════════
\section{I firewall}
% ══════════════════════════════════════════════════════════════

\dfn{Firewall}{
    Un \textbf{firewall} è una combinazione di hardware e software che \textbf{isola} (protegge) una rete da Internet, permettendo o negando il
    passaggio dei pacchetti.
}

\begin{figure}[H]
\centering
\begin{minipage}{0.45\linewidth}\centering
  \includegraphics[width=0.95\linewidth]{cap22/packet_filter.png}
\end{minipage}\hfill
\begin{minipage}{0.51\linewidth}
\begin{itemize}
  \item \textbf{Tutto} il traffico in entrata e in uscita dovrebbe passare per il firewall, e solo il traffico \textbf{autorizzato} (secondo la politica di
        sicurezza locale) può passare.
  \item Il firewall fa da \textbf{unico punto di accesso} alla rete pubblica, ma le grandi organizzazioni possono avere più livelli di firewall o
        firewall distribuiti.
  \item Il firewall stesso deve essere \textbf{immune alle intrusioni}: se viene compromesso, dà un falso senso di sicurezza.
\end{itemize}
\end{minipage}
\caption{Un firewall tra la rete amministrata e la rete pubblica.}
\end{figure}

Vediamo tre approcci: il \textbf{filtraggio tradizionale} dei pacchetti, il \textbf{filtraggio stateful} e l'\textbf{application gateway}.

% ══════════════════════════════════════════════════════════════
\section{Il filtraggio tradizionale dei pacchetti}
% ══════════════════════════════════════════════════════════════

\begin{itemize}
    \item Un \textbf{filtro di pacchetti} è di solito realizzato sul \textbf{gateway}, il router che collega la rete locale all'ISP.
    \item Esamina \textbf{ogni singolo datagramma} e decide se accettarlo o scartarlo in base a regole specificate dall'amministratore. Si possono avere
          regole diverse per i datagrammi in entrata e in uscita, o per le singole interfacce.
\end{itemize}

Le regole si basano tipicamente su:

\begin{itemize}
    \item indirizzo IP sorgente o destinazione;
    \item tipo di protocollo nel campo del datagramma IP (TCP, UDP, ICMP, \dots);
    \item porta sorgente e destinazione;
    \item bit dei flag TCP (SYN, ACK, \dots);
    \item \dots
\end{itemize}

\begin{table}[H]
\centering
\begin{tabular}{p{6cm}p{8cm}}
\toprule
\textbf{Politica} & \textbf{Impostazione del firewall} \\
\midrule
Solo traffico Web in uscita & scartare tutti i segmenti TCP SYN in uscita con porta di destinazione diversa da 80 \\
Niente streaming (non critico) & scartare il traffico UDP in entrata non essenziale \\
Non rispondere ai ping dall'esterno & scartare le risposte ICMP ping in uscita (echo reply) \\
Telnet solo con host fidati & inoltrare i datagrammi Telnet (porta 23) solo da e verso un elenco di indirizzi \\
\bottomrule
\end{tabular}
\caption{Esempi di politiche e relative regole.}
\end{table}

\warning{I filtri non proteggono dallo spoofing}{
    Politiche basate su indirizzi e porte non proteggono dallo \textbf{spoofing}: un pacchetto con indirizzo sorgente falsificato su quello di un host
    fidato passa il filtro.
}

\subsection{Le access control list}

Nei router le regole del firewall sono realizzate con le \textbf{access control list} (ACL). Consideriamo l'ACL dell'interfaccia che collega il
router all'ISP, per un'organizzazione con indirizzi \texttt{222.22/16} che vuole lasciar passare \textbf{solo il traffico Web}:

\begin{table}[H]
\centering
\small
\begin{tabular}{lllllll}
\toprule
\textbf{Azione} & \textbf{IP sorgente} & \textbf{IP destinazione} & \textbf{Protocollo} & \textbf{Porta sorg.} & \textbf{Porta dest.} & \textbf{Flag} \\
\midrule
allow & 222.22/16 & esterno a 222.22/16 & TCP & $>$ 1023 & 80 & qualsiasi \\
allow & esterno a 222.22/16 & 222.22/16 & TCP & 80 & $>$ 1023 & ACK \\
allow & 222.22/16 & esterno a 222.22/16 & UDP & $>$ 1023 & 53 & --- \\
allow & esterno a 222.22/16 & 222.22/16 & UDP & 53 & $>$ 1023 & --- \\
deny & tutto & tutto & tutto & tutto & tutto & tutto \\
\bottomrule
\end{tabular}
\caption{Una ACL che permette solo il traffico Web (le regole si leggono dall'alto: vale la prima che corrisponde).}
\end{table}

\begin{enumerate}
    \item La prima regola lascia \textbf{uscire} qualsiasi pacchetto TCP con porta di destinazione 80.
    \item La seconda lascia \textbf{entrare} qualsiasi pacchetto TCP con porta sorgente 80 e flag \textbf{ACK = 1} (tutti i messaggi Web di risposta portano
          un ACK, mentre un tentativo di aprire una connessione dall'esterno avrebbe SYN senza ACK).
    \item La terza e la quarta insieme permettono ai pacchetti \textbf{DNS} di uscire ed entrare (senza DNS non si navigherebbe).
    \item L'ultima nega tutto il resto.
\end{enumerate}

\subsection{Il limite: il filtraggio è stateless}

\warning{Pacchetti ACK falsi}{
    Il filtraggio tradizionale è \textbf{senza stato}: esamina ogni pacchetto da solo. Per quanto restrittiva, la tabella precedente lascia entrare
    \textbf{qualsiasi} pacchetto dall'esterno con ACK = 1 e porta sorgente 80, anche se non fa parte di nessuna connessione aperta dall'interno.
    Pacchetti di questo tipo sono noti per essere usati negli attacchi DoS. La soluzione ingenua sarebbe bloccare anche i pacchetti TCP ACK, ma
    impedirebbe agli utenti interni di navigare sul Web.
}

% ══════════════════════════════════════════════════════════════
\section{Il filtraggio stateful}
% ══════════════════════════════════════════════════════════════

\dfn{Filtro stateful}{
    I \textbf{filtri stateful} risolvono il problema tenendo traccia di \textbf{tutte le connessioni TCP in corso} in una \textbf{tabella delle
    connessioni}, per capire se il traffico appartiene a una connessione legittima.
    \begin{itemize}
      \item Il firewall riconosce l'\textbf{inizio} di una nuova connessione (la sequenza SYN, SYNACK, ACK del three-way handshake) e la sua
            \textbf{fine} (il pacchetto FIN).
      \item Può anche assumere (in modo conservativo) che la connessione sia finita quando non vede attività per un certo tempo (per esempio 60 secondi).
    \end{itemize}
}

\begin{figure}[H]
    \centering
    \includegraphics[width=0.65\linewidth]{cap22/stateful_table.png}
    \caption{L'ACL di un filtro stateful: la colonna ``check connection'' indica le regole che richiedono di verificare la tabella delle connessioni. In questo caso sono ammessi solo i pacchetti ACK che appartengono a connessioni già stabilite.}
\end{figure}

Una tabella stateful ha di solito una colonna ``\emph{check connection}'' che indica se un pacchetto deve far parte di una connessione già
stabilita: un pacchetto con ACK = 1 in arrivo dall'esterno passa solo se nella tabella c'è una connessione aperta dall'interno con gli stessi
indirizzi e porte.

% ══════════════════════════════════════════════════════════════
\section{L'application gateway}
% ══════════════════════════════════════════════════════════════

I filtri visti finora controllano soprattutto indirizzi IP e porte. Ma può servire permettere o negare funzionalità in base agli \textbf{utenti} o
all'\textbf{applicazione}, per esempio:

\begin{itemize}
    \item solo i tecnici possono usare certi protocolli (come Telnet verso l'esterno), negati agli utenti normali;
    \item all'interno di un protocollo sono permessi solo certi tipi di messaggi o comandi.
\end{itemize}

Compiti del genere vanno oltre le capacità dei filtri tradizionali e stateful: l'identità degli utenti interni è un dato applicativo, e non compare
negli header IP/TCP/UDP.

\dfn{Application gateway}{
    Un \textbf{application gateway} (o \emph{application-level gateway}, ALG) permette di filtrare i pacchetti in base ai \textbf{dati di livello
    applicativo}. È realizzato di solito come un \textbf{server separato} che lavora insieme al firewall.
}

\begin{figure}[H]
\centering
\begin{minipage}{0.4\linewidth}\centering
  \includegraphics[width=0.95\linewidth]{cap22/app_gateway.png}
\end{minipage}\hfill
\begin{minipage}{0.56\linewidth}
\begin{itemize}
  \item L'idea è \textbf{negare tutte le connessioni} di un certo protocollo (per esempio Telnet) \textbf{tranne} quelle da e verso l'application gateway.
  \item Un utente che vuole usare il protocollo ristretto deve prima \textbf{autenticarsi presso il gateway} (nome utente e password).
  \item Il server controlla se l'utente ha il permesso per quel protocollo.
  \item Se sì, tutte le richieste e le risposte vengono inoltrate attraverso il gateway (che fa da \textbf{proxy}).
\end{itemize}
Le reti interne hanno spesso più application gateway, per protocolli diversi (Telnet, HTTP, FTP, posta).
\end{minipage}
\caption{Un gateway per Telnet: sessione host--gateway e sessione gateway--host remoto, con un router che filtra tutto il resto.}
\end{figure}

\warning{Gli svantaggi}{
    \begin{itemize}
      \item Serve un application gateway \textbf{diverso per ogni applicazione}.
      \item C'è una \textbf{penalità di prestazioni} dovuta al proxy (soprattutto con molti utenti).
      \item Il software sugli host deve sapere \textbf{come usare} l'application gateway.
    \end{itemize}
}

% ══════════════════════════════════════════════════════════════
\section{Gli Intrusion Detection System}
% ══════════════════════════════════════════════════════════════

Per rilevare certi tipi di attacco serve un'\textbf{ispezione profonda dei pacchetti} (\emph{deep packet inspection}): controllare, anche nel
contenuto, schemi di attacco noti o traffico sospetto.

\dfn{IDS e IPS}{
    Un \textbf{Intrusion Detection System} (IDS) è un insieme di dispositivi specializzati che \textbf{monitorano la rete} cercando pacchetti sospetti.
    \begin{itemize}
      \item Il suo compito principale è \textbf{riconoscere} i pacchetti potenzialmente malevoli e \textbf{allertare} l'amministratore di rete (che poi
            deciderà cosa fare).
      \item Se può anche \textbf{bloccare} (filtrare) i pacchetti sospetti, si parla di \textbf{Intrusion Prevention System} (IPS).
    \end{itemize}
}

Le organizzazioni usano gli IDS per rilevare un'ampia gamma di attacchi (o di potenziali attacchi):

\begin{itemize}
    \item scansioni della rete o delle porte (per esempio con nmap);
    \item flooding DoS;
    \item worm e virus;
    \item \dots
\end{itemize}

\begin{figure}[H]
    \centering
    \includegraphics[width=0.65\linewidth]{cap22/ids.png}
    \caption{Una rete con filtro, application gateway, sensori IDS e zona demilitarizzata (da Kurose).}
\end{figure}

Il sistema è spesso composto da:

\begin{itemize}
    \item uno o più \textbf{sensori IDS};
    \item un unico \textbf{processore IDS centrale}, che raccoglie e integra le informazioni e invia gli allarmi.
\end{itemize}

Gli IDS possono essere:

\begin{itemize}
    \item \textbf{basati su firme} (\emph{signature-based}): mantengono un ampio database di \textbf{firme di attacco}, cioè insiemi di regole legate a
          un'attività di intrusione (è il tipo più comune). Non riconoscono attacchi nuovi;
    \item \textbf{basati su anomalie} (\emph{anomaly-based}): costruiscono un \textbf{profilo del traffico} normale e segnalano i flussi statisticamente
          insoliti. Possono scoprire attacchi nuovi, ma è difficile distinguere il traffico anomalo da quello semplicemente raro.
\end{itemize}

% ══════════════════════════════════════════════════════════════
\section{La zona demilitarizzata (DMZ)}
% ══════════════════════════════════════════════════════════════

Nelle reti grandi capita spesso di dover avere \textbf{due livelli di sicurezza} per dispositivi diversi: i server che devono comunicare con il mondo
esterno, per esempio, hanno bisogno di un filtraggio più blando.

\dfn{Zona demilitarizzata (DMZ)}{
    Una \textbf{zona demilitarizzata} (DMZ) è una regione a \textbf{bassa sicurezza}, con restrizioni limitate, in cui ospitare i server che devono
    essere raggiungibili dall'esterno (Web, FTP, DNS). Un approccio tipico è mettere la DMZ \textbf{tra due firewall}: uno esterno, meno restrittivo, e
    uno interno, più restrittivo.
}

\begin{figure}[H]
\centering
\begin{tikzpicture}[every node/.style={font=\footnotesize\sffamily}]
  \node[netcloud, minimum width=2.4cm] (i) at (0,0) {Internet};
  \node[draw, fill=netred!30, minimum width=0.5cm, minimum height=1.8cm] (f1) at (2.6,0) {};
  \node[below=2pt of f1, align=center] {firewall\\esterno};
  \node[draw=netorange, thick, dashed, rounded corners, fill=netorange!8, minimum width=3.4cm, minimum height=2.4cm] (dmz) at (5.6,0) {};
  \node[anchor=north, text=netorange, font=\footnotesize\bfseries\sffamily] at (dmz.north) {DMZ};
  \node at (4.8,-0.3) {\icon[0.45cm]{server}}; \node at (5.6,-0.3) {\icon[0.45cm]{server}}; \node at (6.4,-0.3) {\icon[0.45cm]{server}};
  \node[font=\scriptsize\sffamily] at (5.6,-0.95) {Web \quad FTP \quad DNS};
  \node[draw, fill=netred!60, minimum width=0.5cm, minimum height=1.8cm] (f2) at (8.6,0) {};
  \node[below=2pt of f2, align=center] {firewall interno\\+ app. gateway};
  \node[draw=netgreen!60!black, thick, rounded corners, fill=netgreen!8, minimum width=3.2cm, minimum height=2.4cm] (lan) at (11.6,0) {};
  \node[anchor=north, text=netgreen!50!black, font=\footnotesize\bfseries\sffamily] at (lan.north) {rete interna};
  \node at (11,-0.3) {\icon[0.6cm]{pc}}; \node at (12.1,-0.3) {\icon[0.7cm]{laptop}};
  \draw[wired] (i) -- (f1); \draw[wired] (f1) -- (dmz); \draw[wired] (dmz) -- (f2); \draw[wired] (f2) -- (lan);
  \node[text width=4cm, align=center, text=netred] at (5.6,-1.8) {filtraggio blando, sensori IDS};
  \node[text width=4cm, align=center, text=netgreen!50!black] at (11.6,-1.8) {alta sicurezza, sensori IDS};
\end{tikzpicture}
\caption{La DMZ tra due firewall.}
\end{figure}

\begin{itemize}
    \item La regione ad \textbf{alta sicurezza} è protetta da un filtro di pacchetti e da un application gateway, ed è monitorata da sensori IDS.
    \item La regione a \textbf{bassa sicurezza} (la DMZ) è protetta solo dal filtro ed è monitorata dai sensori.
\end{itemize}

\info{Oltre la sicurezza: la privatezza}{
    Crittografia e firewall proteggono i contenuti, ma chi osserva la rete vede comunque \textbf{chi} parla con \textbf{chi} (gli indirizzi IP non sono
    cifrati). Per nascondere anche questo esistono reti di anonimizzazione come \textbf{Tor}, che fa passare il traffico cifrato attraverso più nodi, e
    sistemi come \textbf{Whonix} che instradano tutto il traffico di una macchina virtuale attraverso Tor.
}

% ══════════════════════════════════════════════════════════════
\section{Riepilogo}
% ══════════════════════════════════════════════════════════════

\riepilogo{
\begin{itemize}
    \item \textbf{Autenticazione}: l'IP si falsifica, la password si intercetta, la password cifrata si ripete (\textbf{playback}); la soluzione è il
          \textbf{nonce}, usato una sola volta.
    \item \textbf{SSL/TLS} protegge TCP: \textbf{handshake} (certificato del server, master secret cifrato con la chiave pubblica), \textbf{derivazione}
          di 4 chiavi (cifratura e MAC per ogni direzione), \textbf{trasferimento} a record con MAC e numero di sequenza. HTTPS = HTTP + TLS.
    \item \textbf{Firewall}: filtraggio \textbf{tradizionale} (ACL su indirizzi, porte, flag; stateless), \textbf{stateful} (tabella delle connessioni),
          \textbf{application gateway} (dati applicativi, utenti, proxy).
    \item \textbf{IDS/IPS}: ispezione profonda, basati su firme o su anomalie; la \textbf{DMZ} ospita i server pubblici tra due firewall.
\end{itemize}
}
